\section{GPD 在 $x=\pm\xi$ 附近}\label{sec:yzliu-gpd}

$x=\pm\xi$ 处一个部分子的光锥动量趋于零，“可能的软标度” $2|x\pm\xi|P^z$ 似乎使 LaMET 失效。下面的命题在单圈给出精确的结构，猜想则陈述 Su 等人\cite{Su:2026eta}（预印本）的全阶主张。

\begin{proposition}[准 GPD 单圈匹配核在 $x=\pm\xi$ 附近的结构 \tagR]\label{prop:yzliu-gpd-oneloop}
（$R_{\mathrm{PT}}$：在重整化微扰论中对红外正规化的部分子态严格。）本命题限于单圈阶。考虑\cref{def:lamet-quasi-gpd}的非单态准 GPD（$\Gamma=\gamma^t$，$H$ 型结构，$\overline{\mathrm{MS}}$），外态为动量 $p^z$ 的夸克；以 $p^z$ 重新标度，准变量为 $x'$，斜度为 $0<\xi'<1$（\cref{def:excl-gpd}）。记 $\ell(s)\equiv\ln\frac{4p_z^2s^2}{\mu^2}$。单圈核 $\frac{\alpha_sC_F}{2\pi}G(x',\xi')$ 在 $x'<-\xi'$、$|x'|<\xi'$、$\xi'<x'<1$、$x'>1$ 四个区域中都可写成 % re-derived from the explicit kernel (MD14), script s09 (W26)
\begin{equation}\label{eq:yzliu-gpd-decomp}
G=c_+\,\ell(x'+\xi')+c_-\,\ell(x'-\xi')+c_1\,\ell(1-x')+R,\qquad c_\pm=\frac{|x'\pm\xi'|}{1\pm\xi'}\Big[\frac1{2(1-x')}\pm\frac1{4\xi'}\Big],
\end{equation}
$R$ 为有理函数（另有 $x'=1$ 处的虚修正，即 $\delta(1-x')$ 项与加号处方）。因此可能的软标度 $2|x'\pm\xi'|p^z$ 的对数总乘以 $|x'\pm\xi'|$，在 $x'\to\mp\xi'$ 时线性趋于零。回到原变量：在准 GPD 的匹配公式中（$y$ 为准变量，$x$ 为光锥变量，$|\xi|<x$），这些项为 $\big[\frac1{2(x-y)}\pm\frac1{4\xi}\big]\frac{|y\pm\xi|}{x\pm\xi}\ln\frac{4\bar P_z^2(y\pm\xi)^2}{\mu^2}$；在单圈逆匹配（由准 GPD 求光锥 GPD）中，相应项为 $\big[\frac1{2(x-y)}\mp\frac1{4\xi}\big]\frac{|x\pm\xi|}{y\pm\xi}\ln\frac{4\bar P_z^2(x\pm\xi)^2}{\mu^2}$（$x$ 为光锥变量，$y$ 为准变量，$|\xi|<y$）。（本命题只涉及部分子态。）
\uses{def:lamet-quasi-gpd,def:excl-gpd}
\end{proposition}

\begin{proof}
记 $\ell_\pm=\ln|x'\pm\xi'|$，$\ell_1=\ln|1-x'|$，$L_\mu=\ln\frac{4p_z^2}{\mu^2}$。显式的单圈核\cite[式~(39)--(43)、(46)]{Liu:2019urm}（该文第一作者为 Yu-Sheng Liu）为：$x'<-\xi'$ 时 $G_1=-b_+(\ell_1-\ell_+)-b_-(\ell_1-\ell_-)$，$b_\pm=\frac1{x'-1}\mp\frac{x'}{2\xi'}+\frac{1+x'}{2(1\pm\xi')}$；$|x'|<\xi'$ 时 $G_2=E(L_\mu+2\ell_1+\ell_+-\ell_--1)+D(\ell_--\ell_1)$，$E=\frac{(x'+\xi')(1-x'+2\xi')}{2\xi'(1-x')(1+\xi')}$，$D=\frac{x'+\xi'^2}{\xi'(1-\xi'^2)}$；$\xi'<x'<1$ 时 $G_3=A(L_\mu+\frac12(\ell_++\ell_-)+\ell_1-1)+B(\ell_+-\ell_-)$，$A=\frac{1+x'^2-2\xi'^2}{(1-x')(1-\xi'^2)}$，$B=\frac{x'+\xi'^2}{2\xi'(1-\xi'^2)}$；$x'>1$ 时为 $-G_1$。因 $\ell(s)=L_\mu+2\ln|s|$，形如\cref{eq:yzliu-gpd-decomp}的分解存在当且仅当 $L_\mu$ 的系数等于 $c_++c_-+c_1$，此时 $c_\pm$ 为 $\ell_\pm$ 系数之半、$c_1$ 为 $\ell_1$ 系数之半。

（1）$\xi'<x'<1$：$c_+=\frac A4+\frac B2=\frac{\xi'(1+x'^2-2\xi'^2)+(x'+\xi'^2)(1-x')}{4\xi'(1-x')(1-\xi'^2)}$。分子等于 $(x'+\xi')(1-x'+2\xi')(1-\xi')$（两边展开均为 $x'-x'^2+\xi'+\xi'^2+x'^2\xi'-x'\xi'^2-2\xi'^3$），故 $c_+=\frac{(x'+\xi')(1-x'+2\xi')}{4\xi'(1-x')(1+\xi')}=\frac E2=\frac{x'+\xi'}{1+\xi'}\big[\frac1{2(1-x')}+\frac1{4\xi'}\big]$。$A$ 对 $\xi'$ 为偶、$B$ 为奇，故 $c_-=\frac A4-\frac B2=c_+|_{\xi'\to-\xi'}$，而 $c_++c_-+c_1=\frac A2+\frac A2=A$ 即 $L_\mu$ 的系数。（2）$|x'|<\xi'$：$c_+=\frac E2$，同上；$c_-=\frac{D-E}2$，通分后分子为 $2(x'+\xi'^2)(1-x')-(x'+\xi')(1-x'+2\xi')(1-\xi')=(\xi'-x')(x'+2\xi'-1)(1+\xi')$，故 $c_-=\frac{(\xi'-x')(x'+2\xi'-1)}{4\xi'(1-x')(1-\xi')}=\frac{|x'-\xi'|}{1-\xi'}\big[\frac1{2(1-x')}-\frac1{4\xi'}\big]$；$c_1=E-\frac D2$，三者之和为 $E$，即 $L_\mu$ 的系数。（3）$x'<-\xi'$：$G_1$ 不含 $L_\mu$；通分得 $b_+=-\frac{(x'+\xi')(1-x'+2\xi')}{2\xi'(1-x')(1+\xi')}=-E$，故 $c_+=-\frac E2=\frac{|x'+\xi'|}{1+\xi'}\big[\frac1{2(1-x')}+\frac1{4\xi'}\big]$（此处 $x'+\xi'<0$），同理 $c_-=\frac{b_-}2$ 给出闭式，$c_1=-\frac{b_++b_-}2$ 使三者之和为零。（4）$x'>1$：$G=-G_1$，$c_\pm$ 变号，而 $x'\pm\xi'>0$，闭式仍成立。于是\cref{eq:yzliu-gpd-decomp}在四个区域成立，$c_\pm$ 含因子 $|x'\pm\xi'|$。

换回原变量：$x'=y/x$，$\xi'=\xi/x$，$p^z=x\bar P^z$，并乘以 $1/|x|$，$\frac1x c_+=\frac{|y+\xi|}{x+\xi}\big[\frac1{2(x-y)}+\frac1{4\xi}\big]$，$\ell(x'+\xi')=\ln\frac{4\bar P_z^2(y+\xi)^2}{\mu^2}$。单圈逆核为 $\delta-C^{(1)}$，其中输出变量（光锥 $x$）与积分变量（准 $y$）交换角色，$-\frac1y c_\pm(x/y,\xi/y)$ 即陈述中的形式。

检验：$\xi'\to0$ 时 $G_1$、$G_3$ 趋于准 PDF 的单圈核；$L_\mu$ 的系数在 ERBL 区为 $E$、在 DGLAP 区为 $A$，与由光线算符演化核独立推出的 GPD 演化核一致；ERBL 区 $\ell_+$ 的系数与截断方案的独立计算\cite{Ji:2015qla}一致。
\end{proof}
% Su:2026eta eq. (7) as printed (arXiv v1, rendered page 3; it describes the inverse kernel C of their eq. (6), y quasi, x light-cone)
% has (1/(2(y-x)) -+ 1/(4 xi)) |x+-xi|/(y+-xi): relative to the inverse kernel derived here from Liu:2019urm eqs. (40)-(46)
% (transcription re-checked on the rendered page 7) the 1/(2(y-x)) term has the opposite sign; with the other reading (x quasi)
% the 1/(4 xi) term would disagree instead. Candidate erratum of a preprint, reported, not printed; the vanishing as x -> -+xi
% is unaffected. Scripts s09, s14 (W26).

\paragraph{地位与缺口。}在所述阶为 [R, $R_{\mathrm{PT}}$]。全阶匹配\cref{thm:lamet-quasi-gpd}（[A]）只是指引，不是前提。这是 DVCS 在 $x_B=\xi$ 处的类比：那里 $\xi-x_B$ 的双对数乘以 $(x_B-\xi)$，与实光子相联的软区受幂次压低\cite{JiOsborne1998,CollinsFreund1999,Radyushkin1997}。

\begin{conjecture}[GPD 在 $x=\pm\xi$ 附近的 LaMET 可计算性 \tagC]\label{conj:yzliu-gpd-xi}
在 $x=\pm\xi$ 附近，准 GPD 的匹配不受与 $(x\pm\xi)P^z$ 相联的软标度的领头扭度贡献；条件 $2|x\pm\xi|P^z\gg\Lambda$ 可代之以 $2(|x|+|\xi|)P^z\gg\Lambda$ 与 $2|1\pm x|P^z\gg\Lambda$；以完整单圈匹配提取的 GPD 在 $x=\pm\xi$ 处连续而导数不连续（取正规子次序 $\epsilon_2\ll|\epsilon_1|\to0$；其他次序发散）\cite{Su:2026eta}。非微扰地位：\cref{asm:intro-bridge}、\cref{asm:intro-composite}（MA0、MA1、MA2）。
\uses{prop:yzliu-gpd-oneloop,thm:lamet-quasi-gpd}
\end{conjecture}

\paragraph{地位与缺口。}[C]，开放。文献\cite{Su:2026eta}（预印本）对四组格点数据作了重新分析，认为此前看到的不连续“may arise from regulator-related numerical artifacts”；单圈的依据是\cref{prop:yzliu-gpd-oneloop}。另一方面，重整子分析预言 $x=\pm\xi$ 处的非微扰修正只受硬标度一次幂压低，来自软动量，与 Wilson 线紫外重整子无关\cite{Braun:2024snf}；运动学扭度 3 贡献在 $x=\pm\xi$ 处对数增强且不连续\cite{Braun:2023alc}。两种立场并不严格矛盾（没有领头扭度的软项，与增强到 $1/P^z$ 的幂次修正），两方也未互相批评。证明需要对准 GPD 在 $x=\pm\xi$ 处的软夸克（零模）区作全阶幂次计数，类似\cite{CollinsFreund1999}。
